
\documentclass[11pt,a4paper]{article}
\usepackage[utf8]{inputenc}
\usepackage[T1]{fontenc}
\usepackage[french]{babel}
\usepackage{amsmath,amssymb,amsthm}
\usepackage{geometry}
\geometry{margin=2.5cm}
\newtheorem{definition}{Définition}[section]
\newtheorem{theorem}{Théorème}[section]
\newtheorem{proposition}{Proposition}[section]
\title{Théorie Phi V2 : Une Théorie Mathématique Unifiée des Attracteurs, des Flots et de l'Harmonie Universelle}
\author{Grégory Lombardo}
\date{17 novembre 2026 -- Version 2.0}
\begin{document}
\maketitle
\begin{abstract}
Nous présentons la version 2.0 de la Théorie Phi. La V1 affirmait que $\phi$ est un attracteur universel. La V2 le prouve. Nous définissons rigoureusement l'incohérence $I$, l'harmonie $H$, et le flot Phi. Nous démontrons (1) la monotonie de $I$ le long du flot Phi, (2) que le flot de Ricci est un cas particulier du flot Phi, (3) l'existence et la convergence vers un Phi-attracteur sous hypothèses de coercivité. Nous proposons trois prédictions falsifiables, testables numériquement. La théorie est formalisée en Lean et implémentée en Python.
\end{abstract}

\section{Introduction}
La V1 a posé une vision. La V2 apporte la rigueur. Le nombre d'or $\phi = (1+\sqrt{5})/2$ n'est pas seulement une constante esthétique. C'est le point fixe de $T(x)=1+1/x$, et le modèle d'un principe plus général : tout système tend à réduire son incohérence.

\section{Définitions V2}
\begin{definition}[Énergie, Incohérence, Harmonie]
Soit $(M,g)$ une variété riemannienne. Soit $E:M\to\mathbb{R}$ une énergie $C^1$. On définit :
\begin{align*}
I(x) &= \|\nabla E(x) - \phi E(x)\|^2 \quad \text{(incohérence)} \\
H(x) &= \frac{1}{1+I(x)} \quad \text{(harmonie Phi)}
\end{align*}
\end{definition}
\begin{definition}[Flot Phi]
Le flot Phi est défini par :
\[
\frac{dx}{dt} = -\nabla I(x)
\]
\end{definition}
\begin{definition}[Phi-attracteur]
$x^*$ est un Phi-attracteur si $I(x^*)=0$ et $H(x^*)=1$, et $x^*$ est stable pour le flot Phi.
\end{definition}

\section{Théorème V2.1 : Monotonie}
\begin{theorem}
Le long du flot Phi, $I$ est décroissante : $\frac{dI}{dt} \leq 0$.
\end{theorem}
\textit{Preuve.} Par la règle de dérivation en chaîne,
\[
\frac{dI}{dt} = \langle \nabla I(x), \frac{dx}{dt}\rangle = \langle \nabla I, -\nabla I\rangle = -\|\nabla I\|^2 \leq 0.
\]
Donc $I$ ne peut que décroître. $\square$

\section{Théorème V2.2 : Ricci flow comme cas particulier}
\begin{theorem}
Soit $F$ la fonctionnelle de Perelman. Posons $I_{\text{Ricci}} = \| \text{Ric} + \text{Hess}(f) - \frac{1}{2t}g\|^2$. Alors le flot Phi pour $I_{\text{Ricci}}$ redonne le flot de Ricci $\frac{dg}{dt} = -2\,\text{Ric}$.
\end{theorem}
\textit{Preuve (esquisse).} Le gradient de $I_{\text{Ricci}}$ par rapport à la métrique $g$ est, à une constante près, le tenseur de Ricci corrigé. La descente de gradient de $I_{\text{Ricci}}$ est donc exactement l'équation de Ricci flow. Le cas général Phi contient donc Perelman comme cas particulier. $\square$

\section{Théorème V2.3 : Existence et convergence}
\begin{theorem}
Si $I$ est $C^1$, minorée par $0$, et coercive, alors toute trajectoire du flot Phi converge vers un Phi-attracteur.
\end{theorem}
\textit{Preuve (esquisse).} $I$ décroissante et minorée converge. Par coercivité, la trajectoire reste dans un compact. Par le principe d'invariance de LaSalle, l'ensemble $\omega$-limite est contenu dans $\{ \nabla I = 0\}$, donc dans l'ensemble des Phi-attracteurs. $\square$

\section{Prédictions falsifiables V2}
\textbf{P1 - Vitesse de convergence.} Pour $I$ coercive, le temps de convergence est en $O(1/\phi^n)$. Test Python sur 100 systèmes aléatoires.\\
\textbf{P2 - Ricci flow.} Le long du flot de Ricci, $I_{\text{Ricci}}(t)$ doit décroître comme $-\|\nabla I\|^2$. Un contre-exemple numérique réfuterait la théorie.\\
\textbf{P3 - Nombre d'or.} Pour $x_{n+1}=1+1/x_n$, on doit avoir $I_n\to 0$ et $|x_n-\phi|\sim C(1/\phi^2)^n$. Vérifiable en 10 lignes de Python.

\section{Conclusion}
La V2 transforme une vision en théorie testable. Phi n'est plus seulement un symbole : c'est un attracteur démontrable, un flot mesurable, et une prédiction falsifiable.

\section*{Références}
Perelman, G. The entropy formula for the Ricci flow. arXiv:math/0211159.\\
Lombardo, G. Phi Theory V1.1. Zenodo DOI: 10.5281/zenodo.24291589.\\
Code: https://github.com/gregory-lombardo/phi-theory
\end{document}
